\chapter{Functional Testing}
\label{sec:functionalTesting}
This chapter details the testing which was done on the system's primary function: velocity measurement. Two distinct test setups are used, which are described in this section - results were then collected for these test setups and are presented and discussed with respect to the goals of the system and the overarching project. 

\section{C: Hard-target Speed Measurement}
\subsection{Experimental Goals}
The overarching goal of the system is to measure wind-speed via the Doppler shift, however, owing to the very weak signal return in air and complexities in setting up appropriate tests, this is a challenging test to perform. For this reason it was decided that before wind-speed measurement testing is performed, hard-target speed measurement would first be tested. This uses the same working principle as wind-speed measurement, but with a target which provides a far higher signal return. Thus, this experiment serves to first validate the working principle of the system to measure velocity via the optical Doppler shift, before wind-speed measurement testing is performed. This allows for any significant setup or procedure problems in the system to be identified in an environment which is easier to control, and which is more reliable for measurement (owing to the larger return signal, and thus higher SNR). 

\subsection{Experimental Methodology}
In order to generate a constant and precise hard-target velocity to be measured, a NEMA-17 stepper motor is used with a disc attached to the motor shaft. Stepper motors are capable of rotating at very constant and accurate RPMs. If the motor axis is tilted at an angle with respect to the laser beam, there will be a component of the disc's angular velocity in the direction of the laser beam. If the beam is carefully aligned with a known point on the disc's radius, the instantaneous velocity of that point in the direction of the beam, $v_{los}$, can thus be calculated. A design was created and 3D printed, which aligns the collimator lens with a spot on the spinning disc's radius. \autoref{fig:CRig} shows the setup used:
\begin{figure}[H]
	\centering
	\includegraphics[width=1\linewidth]{../../../../../Downloads/fig_rig_side3d}
	\caption{Test setup used for hard-target speed measurement}
	\label{fig:CRig}
\end{figure}
As shown in the figure, an angle of 45 degrees is used between the collimator and the motor axis, and the beam is aimed approximately $d = 15mm$ from the centre of the disc. The line of sight velocity of the beam is thus: $v_{los}=sin(45)2\pi N \times d$, where $N$ is the RPM of the motor, and $d$ is the radius from the centre of the disc to the beam spot. While care was taken in the design to precisely align the collimator with the intended point at $d=15mm$, any error in this alignment will translate directly into a measured velocity which is different from the calculated velocity. Furthermore, the beam spot has a finite size which has a diameter of approximately 7mm (from the F810-APC datasheet \cite{f810 apc datasheet}) - thus the measurement will also have contributions from other points on the disc radius which move at a different angular velocity to the intended target point. These contributions are likely to appear as a broadening of the Doppler peak. 

It must also be noted that the PLA plastic used reflects very little light at 1550nm, and thus standard white paper was attached to the spinning disc to increase the reflectivity. 

The rest of the system is set up as in \autoref{fig:finalDesign}, with all components fully connected. The laser is driven at a constant 48mA.

The motor RPM is varied across measurements in order to test the system's ability to resolve various different speeds. The same data acquisition and processing chain as for experiments A and B are used here, with 100 oscilloscope records being obtained for a given motor RPM, their Welch PSDs computed, and averaged to obtain a result. 

\subsection{Results and Discussion}

\begin{figure}[H]
	\centering
	\includegraphics[width=1\linewidth]{"../Data Acquisition/figures/psd_rpm_overlay"}
	\caption{C: PSD of hard-target speed measurements, at different motor RPMs}
	\label{fig:CPSD}
\end{figure}

\begin{figure}
	\centering
	\includegraphics[width=1\linewidth]{"../Data Acquisition/figures/peak_vs_rpm"}
	\caption{C: Doppler peak frequency (from centroid) vs target motor RPM}
	\label{fig:CentroidVSRPM}
\end{figure}

\begin{table}[H]
	\centering
	\small
	\setlength{\tabcolsep}{4pt}
	\caption{Doppler peak and local noise floor for each target speed (capture~1 of each speed,
		Welch PSD with $\Delta f = \SI{10}{kHz}$). $f_D$ is the power-weighted centroid of
		the excess PSD over the stationary-target reference, and
		$v_{\mathrm{LOS}} = f_D / (\SI{1.29}{MHz.s/m})$. The PSD at $f_D$ is that of the bin
		nearest the centroid; the noise statistics are taken in a band outside
		$3\times$FWHM around the centroid, out to \SI{2}{MHz} either side. All PSD values are
		in the same units, so the PSD at $f_D$ can be compared directly with the detection
		threshold (mean $+\,5\sigma$), which it exceeds at every speed. The excess power is the PSD above the local
		noise, integrated across the peak.}
	\label{tab:doppler_vs_rpm}
	\smallskip

	\begin{tabular}{
			|S[table-format=4.0]   % RPM
			|S[table-format=1.3]   % f_D (centroid)
			|S[table-format=1.3]   % v_LOS
			|S[table-format=1.2]   % mean noise
			|S[table-format=1.2]   % mean + 5 sigma
			|S[table-format=2.1]   % PSD at f_D
			|S[table-format=2.1]   % PSD/mean
			|S[table-format=1.2]|  % integrated excess power
		}
		\hline
		\multicolumn{1}{|c|}{\textbf{Speed}} & \multicolumn{1}{c|}{\boldmath$f_D$} & \multicolumn{1}{c|}{\boldmath$v_{\mathrm{LOS}}$}
		& \multicolumn{3}{c|}{\textbf{PSD} (\num{e-11}\,\si{V^2/Hz})} & \multicolumn{1}{c|}{\textbf{PSD/mean}} & \multicolumn{1}{c|}{\textbf{Excess power}} \\ \cline{4-6}
		\multicolumn{1}{|c|}{(RPM)} & \multicolumn{1}{c|}{(\si{MHz})} & \multicolumn{1}{c|}{(\si{m/s})}
		& \multicolumn{1}{c|}{Mean noise} & \multicolumn{1}{c|}{Mean $+\,5\sigma$}
		& \multicolumn{1}{c|}{\textbf{At \boldmath$f_D$}} & \multicolumn{1}{c|}{(\si{dB})} & \multicolumn{1}{c|}{(\num{e-5}\,\si{V^2})} \\ \hline
		50 & 0.079 & 0.061 & {7.01\textsuperscript{\dag}} & 9.35 & 50.7 &  8.6 & 4.09 \\ \hline
		100 & 0.164 & 0.127 & {2.56\textsuperscript{\dag}} & 3.95 & 75.7 & 14.7 & 4.00 \\ \hline
		200 & 0.330 & 0.256 & {1.43\textsuperscript{\dag}} & 2.24 & 78.6 & 17.4 & 4.35 \\ \hline
		600 & 0.996 & 0.772 & 4.92 & 7.86 & 26.1 &  7.2 & 3.76 \\ \hline
		800 & 1.331 & 1.032 & 5.61 & 8.99 & 25.2 &  6.5 & 3.91 \\ \hline
		1200 & 1.985 & 1.539 & 4.07 & 6.38 & 14.5 &  5.5 & 3.63 \\ \hline
		2000 & 3.314 & 2.569 & 2.92 & 4.06 & 10.4 &  5.5 & 3.93 \\ \hline
	\end{tabular}
	
	\vspace{2pt}
	\raggedright\footnotesize
	\textsuperscript{\dag}\,Peak too close to DC for a noise band below it; the floor is taken
	from the high-frequency side only and is likely a slight underestimate.
\end{table}

\section{C: Wind-speed Measurement}
\subsection{Experimental Goals}
Now that the noise analysis of the system is complete, and the speed-measurement principle has been shown to be functional for a hard-target, the system can finally be tested for seeded air wind-speed measurement. The primary subject of investigation for this section is the sensitivity of the system with respect to its noise and the different levels of seeding in the air, as well as its accuracy at different tested speeds. As discussed previously, the range of the system can not be tested in a wind-tunnel environment, and thus this is out of scope for this project. 
\subsection{Experimental Setup}
% Wind-tunnel test set-up (side view). Coordinates are in mm, scaled 1 mm -> 0.25 mm.
% Needs: \usepackage{tikz} \usetikzlibrary{arrows.meta,patterns,decorations.pathmorphing,calc}
\begin{tikzpicture}[x=0.25mm, y=0.25mm, font=\small,
	blk/.style={draw, thick, rounded corners=2pt, align=center, fill=gray!8, inner sep=5pt},
	flow/.style={-{Stealth[length=6pt]}, thick, blue!60!black},
	lbl/.style={font=\footnotesize, align=center}]
	
	% ---------------- wind-tunnel test section ----------------
	\fill[pattern=north east lines, pattern color=brown!70] (0,250) rectangle (500,262);
	\fill[pattern=north east lines, pattern color=brown!70] (0,-12) rectangle (500,0);
	\draw[very thick] (0,250) -- (500,250);
	\draw[very thick] (0,0) -- (500,0);
	\draw[dashed, gray] (0,0) -- (0,250);
	\draw[dashed, gray] (500,0) -- (500,250);
	\node[lbl, anchor=south] at (250,264) {MDF ceiling panel (beam terminates here)};
	\node[lbl, anchor=north west, gray!40!black] at (3,246) {inlet};
	\node[lbl, anchor=north east, gray!40!black] at (497,246) {outlet};
	\node[lbl, gray!40!black] at (395,70) {Wind-tunnel test section\\[-1pt]$500 \times 250 \times 250$~mm};
	
	% air-flow arrows
	\foreach \y in {95,150,205} \draw[flow] (12,\y) -- (72,\y);
	\node[lbl, blue!60!black, anchor=west] at (78,150) {air flow};
	
	% ---------------- seeding ----------------
	\node[blk, fill=gray!15] (seed) at (-120,165) {Seeding\\smoke/fog};
	\draw[-{Stealth[length=6pt]}, thick, decorate,
	decoration={snake, amplitude=2pt, segment length=8pt, post length=6pt}]
	(seed.east) -- (-4,165);
	
	% ---------------- wedge and collimator ----------------
	% printed PLA wedge (side profile), lens centre at (105,26)
	\filldraw[fill=orange!35, draw=black, thick]
	(43.2,0) -- (43.2,1.0) -- (96.4,38.3) -- (123.2,0) -- cycle;
	% collimator body along the beam axis (35 deg), lens face at (105,26)
	\begin{scope}[shift={(105,26)}, rotate=35]
		\filldraw[fill=gray!60, draw=black] (-46.8,-6) rectangle (0,6);
		\filldraw[fill=gray!30, draw=black] (-62,-3) rectangle (-46.8,3);   % FC/APC connector
	\end{scope}
	\draw[thin] (200,8) -- (118,8);
	\node[lbl, anchor=west] at (202,8) {PLA wedge + collimator};
	
	% ---------------- beam ----------------
	\draw[red!80!black, very thick] (105,26) -- (425.3,250);
	\fill[red!80!black] (425.3,250) circle (3.5);
	\path (105,26) -- (425.3,250) node[pos=0.62, sloped, above, lbl, red!70!black] {beam, $\lambda = \SI{1550}{nm}$};
	
	% angle to the flow
	\draw[dashed, gray] (105,26) -- (190,26);
	\draw[thick] (160,26) arc[start angle=0, end angle=35, radius=55];
	\node[lbl] at (180,44) {$35^\circ$};
	
	% ---------------- fibre and lidar system ----------------
	\node[blk, minimum width=150, minimum height=46] (sys) at (105,-120)
	{Lidar system\\[-1pt]{\footnotesize laser, circulator, LO, balanced detector, DAQ}};
	\coordinate (rear) at ($(105,26)+(215:62)$);
	\draw[orange!80!black, very thick] (rear) -- (rear |- sys.north);
	\node[lbl, orange!50!black, anchor=west] at (60,-58) {single-mode fibre (FC/APC)};
\end{tikzpicture}

